Покупатель прикладывает карту или нажимает кнопку на странице оплаты,
терминал печатает чек или приложение показывает успешный статус --
и~кажется, что произошло одно событие, в~котором деньги \enquote{перешли}
от плательщика к~получателю. На самом деле в~платеже задействованы
по меньшей мере три разных сущности: сообщение о~намерении совершить
операцию, денежное обязательство между участниками и собственно
расчёт, который может произойти позже, другим каналом и~по другим правилам.

Если не разделить эти три сущности, в~обращении с~базовыми понятиями книги возникнет путаница.
\textit{Авторизация} становится \textit{переводом} денег. \textit{Клиринг}
превращается в~подтверждение того, что и~так случилось,
а регуляторные требования -- в~посторонний нарост на техническом потоке.

Платёжная система связывает их между собой. Она фиксирует,
кто и~когда меняет запись в~своём учёте и~в~какой момент операция
становится необратимой. До этого момента её ещё можно отозвать,
оспорить, скорректировать.

Часть начинается с~трёх базовых вопросов: что вообще считается
деньгами в~банковской инфраструктуре, кто владеет обязательствами
и как одно учреждение узнаёт, что другое должно изменить состояние
счёта. Дальше -- формы денег и~различие между банковским переводом
и карточной операцией. После этого -- сама карта, её идентификаторы
и признаки подлинности; без этого PAN, ответ авторизации
и фактический расчёт сливаются в~одну плоскость, хотя относятся к~разным.

После этой подготовки темы разворачиваются в~трёх плоскостях, к~которым возвращает каждая глава части.
\begin{itemize}
  \item Жизненный цикл и~базовые этапы карточной операции -- отклонения от них рассматриваются в~последующих частях.
  \item Контроль безопасности -- как система убеждается, что операция допустима, карта подлинна, держатель легитимен.
  \item Юридические роли и~то, почему техническая архитектура оказывается одновременно правовой.
\end{itemize}
